\documentclass[12pt, a4paper]{article}

% ─── Encoding e lingua ────────────────────────────────────────────────────────
\usepackage[T1]{fontenc}
\usepackage[utf8]{inputenc}
\usepackage[italian]{babel}

% ─── Layout ───────────────────────────────────────────────────────────────────
\usepackage[top=2.5cm, bottom=2.5cm, left=2.8cm, right=2.8cm, headheight=15pt]{geometry}
\usepackage{parskip}
\setlength{\parindent}{0pt}

% ─── Font ─────────────────────────────────────────────────────────────────────
\usepackage{lmodern}

% ─── Grafica e colori ─────────────────────────────────────────────────────────
\usepackage{graphicx}
\usepackage[table]{xcolor}
\definecolor{headerbg}{RGB}{30, 80, 140}
\definecolor{rowalt}{RGB}{235, 242, 252}
\definecolor{best}{RGB}{198, 239, 206}

% ─── Tabelle ──────────────────────────────────────────────────────────────────
\usepackage{booktabs}
\usepackage{array}
\usepackage{float}
\newcolumntype{C}[1]{>{\centering\arraybackslash}p{#1}}
\newcolumntype{L}[1]{>{\raggedright\arraybackslash}p{#1}}

% ─── Figure affiancate ────────────────────────────────────────────────────────
\usepackage{subcaption}

% ─── Link e metadati ──────────────────────────────────────────────────────────
\usepackage[hidelinks]{hyperref}

% ─── Matematica ───────────────────────────────────────────────────────────────
\usepackage{amsmath}

% ─── Sezioni compatte ─────────────────────────────────────────────────────────
\usepackage{titlesec}
\titlespacing*{\section}{0pt}{14pt}{6pt}
\titlespacing*{\subsection}{0pt}{10pt}{4pt}

% ─── Header / Footer ──────────────────────────────────────────────────────────
\usepackage{fancyhdr}
\pagestyle{fancy}
\fancyhf{}
\fancyhead[L]{\small HYDRAS --- Report FCM}
\fancyhead[R]{}
\fancyfoot[C]{\thepage}
\renewcommand{\headrulewidth}{0.4pt}

% ══════════════════════════════════════════════════════════════════════════════
\begin{document}

% ─── Titolo ───────────────────────────────────────────────────────────────────
\begin{center}
    {\LARGE\bfseries HYDRAS --- Report FCM}\\[6pt]
    {\large Field Climbing Method come Baseline di Confronto con RL}
\end{center}
\vspace{0.3cm}
\hrule
\vspace{0.5cm}

% ══════════════════════════════════════════════════════════════════════════════
\section{Motivazione}

Per valutare il contributo dell'apprendimento per rinforzo nel problema di source
seeking, \`e stato implementato e testato il \textbf{Field Climbing Method} (FCM)
come baseline algoritmica.
FCM rappresenta l'approccio gradient-based classico al problema: stima il gradiente
del campo di concentrazione a partire dalle letture locali dei sensori e si muove
nella direzione di salita pi\`u ripida.
A differenza del modello RL, non richiede training e non mantiene memoria globale
dell'episodio: decide l'azione ad ogni step basandosi esclusivamente sulle
osservazioni correnti.

% ══════════════════════════════════════════════════════════════════════════════
\section{Algoritmo Field Climbing Method}

\subsection{Stima del Gradiente per Minimi Quadrati}

L'agente legge la concentrazione in \textbf{9 punti}: la posizione corrente
$\mathbf{p}$ (concentrazione centrale $C_0$) e 8 sensori direzionali che
campionano $C(\mathbf{p} + \boldsymbol{\delta}_i)$ con offset
$\boldsymbol{\delta}_i = r_s\,\hat{d}_i$, dove $r_s$ \`e il \emph{sensor range}
e $\hat{d}_i$ \`e uno degli 8 versori cardinali e diagonali
(N, S, E, O, NE, SE, NO, SO) --- gli stessi delle azioni discrete (sez.~2.2).

Assumendo un'approssimazione lineare locale del campo (Taylor al primo ordine):
\begin{equation}
    C(\mathbf{p} + \boldsymbol{\delta}_i) \approx C(\mathbf{p}) + \nabla C \cdot \boldsymbol{\delta}_i
\end{equation}

Sottraendo $C_0$ da ciascuna lettura e normalizzando per $r_s$, l'equazione~(1)
diventa $\hat{d}_i \cdot \nabla C \approx b_i$ per ogni sensore $i$, con:
\begin{equation}
    b_i = \frac{C_i - C_0}{r_s}, \qquad C_i \equiv C(\mathbf{p}+\boldsymbol{\delta}_i)
\end{equation}

Raccogliendo gli 8 vincoli in forma matriciale ($\mathbf{D} \in \mathbb{R}^{8 \times 2}$,
$D_{ij} = (\hat{d}_i)_j$) si ottiene il sistema sovradeterminato
$\mathbf{D}\,\nabla C \approx \mathbf{b}$, risolto per minimi quadrati con
\texttt{numpy.linalg.lstsq}. La disposizione simmetrica delle 8 direzioni
(per ogni $\hat{d}_i$ \`e presente anche $-\hat{d}_i$) implica
$\sum_i (\hat{d}_i)_x^2 = \sum_i (\hat{d}_i)_y^2 = 4$ e
$\sum_i (\hat{d}_i)_x(\hat{d}_i)_y = 0$, ovvero
$\mathbf{D}^\top\mathbf{D} = 4\,\mathbf{I}_2$.
Il sistema \`e perfettamente condizionato e la stima ha forma chiusa:
\begin{equation}
    \hat{\nabla C} = \frac{1}{4}\,\mathbf{D}^\top\mathbf{b}
    \;=\; \frac{1}{4}\sum_{i=1}^{8} \hat{d}_i\,b_i
\end{equation}

\subsection{Selezione dell'Azione}

Le 8 azioni discrete corrispondono esattamente alle 8 direzioni $\hat{d}_i$ usate
dai sensori in sez.~2.1, con passo fisso pari a
$v_{\max} \cdot \Delta t = 1\,\text{m/s} \times 10\,\text{s} = 10\,\text{m}$.

Avendo già il gradiente stimato $\hat{\nabla C}$, si sceglie la direzione
discreta pi\`u allineata ad esso:
\begin{equation}
    a^* = \arg\max_{i \in \{1,\ldots,8\}} \; \hat{d}_i \cdot \hat{\nabla C}
\end{equation}
L'agente si sposta quindi di $10\,\text{m}$ nella direzione $\hat{d}_{a^*}$.
Se $\|\hat{\nabla C}\| < 10^{-12}$ (campo localmente piatto),
l'azione viene scelta uniformemente a caso tra le direzioni valide.
In entrambi i casi, le direzioni che porterebbero l'agente fuori dal dominio
(collisione con la costa) vengono preventivamente escluse tramite action masking.

% ══════════════════════════════════════════════════════════════════════════════
\section{Dataset e Ambiente di Simulazione}

Il setup sperimentale \`e identico a quello usato per il modello RL.
Il dataset \`e composto da simulazioni MIKE21 su un dominio marino di
$3 \times 2{,}5\,\text{km}$ (baia di Cecina, UTM32N), risoluzione $10\,\text{m}$,
passo temporale $10\,\text{s}$, con 132 sorgenti $\times$ 4 scenari vento (V0--V3).

\begin{itemize}
    \item \textbf{Sorgenti di inferenza:} SRC107--SRC132 (26 sorgenti)
    \item \textbf{Chunk temporali:} Q1/4, Q1/2, Q3/4
    \item \textbf{Episodi per scenario:} 2
    \item \textbf{Totale episodi per run:} $26 \times 4 \times 3 \times 2 = 624$
    \item \textbf{Durata massima episodio:} 1080 step ($= 3\,\text{h}$ simulate)
    \item \textbf{Soglia di successo:} distanza dalla sorgente $\leq 50\,\text{m}$
    \item \textbf{Distanza iniziale media:} $\approx 800\,\text{m}$
\end{itemize}

% ══════════════════════════════════════════════════════════════════════════════
\section{Sweep del Sensor Range}

Il parametro critico di FCM \`e il \textbf{sensor range} $r_s$: la distanza a cui
vengono campionati i sensori direzionali.
Un range troppo piccolo rende i sensori insensibili al gradiente (differenze di
concentrazione trascurabili su scala locale); un range troppo grande campiona zone
distanti dalla posizione corrente, riducendo la validit\`a dell'approssimazione lineare
e aumentando il tempo medio di convergenza verso la sorgente.

Sono stati testati cinque valori: $r_s \in \{20, 50, 100, 150, 200\}\,\text{m}$,
con 2 episodi per scenario (624 episodi totali per run).

\begin{table}[htbp]
\centering
\renewcommand{\arraystretch}{1.3}
\small
\begin{tabular}{C{2cm} C{2.2cm} C{2.2cm} C{2cm} C{2cm} C{2cm} C{2cm}}
\toprule
\rowcolor{headerbg}
\textcolor{white}{\textbf{$r_s$}} &
\textcolor{white}{\textbf{SR globale}} &
\textcolor{white}{\textbf{Step medi (succ.)}} &
\textcolor{white}{\textbf{V0}} &
\textcolor{white}{\textbf{V1}} &
\textcolor{white}{\textbf{V2}} &
\textcolor{white}{\textbf{V3}} \\
\midrule
20\,m  & 19{,}4\% & 274{,}5 ($\sim$45{,}7\,min) & 17{,}9\% & 10{,}3\% & 19{,}2\% & 30{,}1\% \\
\rowcolor{best}
\textbf{50\,m}  & \textbf{32{,}9\%} & \textbf{182{,}7 ($\sim$30{,}5\,min)} & \textbf{39{,}7\%} & \textbf{16{,}0\%} & \textbf{23{,}7\%} & \textbf{51{,}9\%} \\
100\,m & 18{,}8\% & 411{,}6 ($\sim$68{,}6\,min) & 28{,}2\% & 8{,}3\%  & 16{,}0\% & 22{,}4\% \\
150\,m & 13{,}6\% & 546{,}4 ($\sim$91{,}1\,min) & 16{,}0\% & 5{,}1\%  & 12{,}2\% & 21{,}2\% \\
200\,m & 13{,}8\% & 398{,}7 ($\sim$66{,}4\,min) & 6{,}4\%  & 6{,}4\%  & 13{,}5\% & 28{,}8\% \\
\bottomrule
\end{tabular}
\caption{Risultati dello sweep su sensor range. In verde il valore ottimale ($r_s = 50\,\text{m}$).}
\end{table}

\begin{table}[htbp]
\centering
\renewcommand{\arraystretch}{1.3}
\small
\begin{tabular}{C{2cm} C{3cm} C{3cm} C{3cm}}
\toprule
\rowcolor{headerbg}
\textcolor{white}{\textbf{$r_s$}} &
\textcolor{white}{\textbf{Q1/4}} &
\textcolor{white}{\textbf{Q1/2}} &
\textcolor{white}{\textbf{Q3/4}} \\
\midrule
20\,m  & 41{,}3\% & 1{,}0\%  & 15{,}9\% \\
\rowcolor{best}
\textbf{50\,m}  & \textbf{50{,}5\%} & \textbf{11{,}5\%} & \textbf{36{,}5\%} \\
100\,m & 28{,}8\% & 1{,}4\%  & 26{,}0\% \\
150\,m & 16{,}8\% & 1{,}0\%  & 23{,}1\% \\
200\,m & 11{,}1\% & 7{,}7\%  & 22{,}6\% \\
\bottomrule
\end{tabular}
\caption{Success rate per chunk temporale al variare del sensor range.}
\end{table}

Il valore ottimale risulta $r_s = 50\,\text{m}$, con un success rate globale del
\textbf{32{,}9\%} e un tempo medio al successo di $\sim$30 minuti.
Per valori superiori il success rate crolla significativamente rispetto al valore
ottimale, e il tempo medio al successo aumenta, a causa dell'approssimazione lineare
sempre meno valida su scale spaziali grandi rispetto alla struttura del plume.

% ══════════════════════════════════════════════════════════════════════════════
\section{Analisi dei Risultati --- $r_s = 50\,\text{m}$}

I grafici seguenti analizzano il comportamento di FCM con il sensor range ottimale
su 624 episodi.

\begin{figure}[H]
  \centering
  \includegraphics[width=0.88\linewidth]{../../evaluations/evaluations_FCM/fcm_pure_50m/analysis/plot_initial_dist.png}
  \caption{Distribuzione delle distanze iniziali di spawn --- FCM $r_s = 50\,\text{m}$.}
\end{figure}

\vspace{1.2cm}

\begin{figure}[H]
  \centering
  \includegraphics[width=0.88\linewidth]{../../evaluations/evaluations_FCM/fcm_pure_50m/analysis/plot_sr_analysis.png}
  \caption{Success rate per scenario di vento e chunk temporale --- FCM $r_s = 50\,\text{m}$.}
\end{figure}

\begin{figure}[H]
  \centering
  \includegraphics[width=0.88\linewidth]{../../evaluations/evaluations_FCM/fcm_pure_50m/analysis/plot_success_time_dist.png}
  \caption{Distribuzione del tempo al successo --- FCM $r_s = 50\,\text{m}$.}
\end{figure}

\vspace{1.2cm}

\begin{figure}[H]
  \centering
  \includegraphics[width=0.88\linewidth]{../../evaluations/evaluations_FCM/fcm_pure_50m/analysis/plot_distance_over_time.png}
  \caption{Distanza media dalla sorgente nel tempo --- FCM $r_s = 50\,\text{m}$.}
\end{figure}

% ══════════════════════════════════════════════════════════════════════════════
\subsection{Casi rappresentativi}

\subsubsection*{Episodi di successo}

\begin{figure}[H]
  \centering
  \includegraphics[width=0.62\linewidth]{../../evaluations/evaluations_FCM/fcm_pure_50m/SRC123/V0/ep02_chunk1_trajectory.png}
  \caption{Successo --- SRC123, V0, Q1/2: convergenza in 137 step ($\sim$23 min) da 640\,m.
           Gradiente ben definito lungo tutto l'episodio; risalita direzionale del plume.}
\end{figure}

\vspace{1.2cm}

\begin{figure}[H]
  \centering
  \includegraphics[width=0.62\linewidth]{../../evaluations/evaluations_FCM/fcm_pure_50m/SRC118/V0/ep02_chunk2_trajectory.png}
  \caption{Successo --- SRC118, V0, Q3/4: convergenza in 252 step ($\sim$42 min) da 651\,m.
           Alcune deviazioni intermedie ma tendenza direzionale sostenuta verso la sorgente.}
\end{figure}

\subsubsection*{Episodi di fallimento}

\begin{figure}[H]
  \centering
  \includegraphics[width=0.62\linewidth]{../../evaluations/evaluations_FCM/fcm_pure_50m/SRC109/V1/ep01_chunk1_trajectory.png}
  \caption{Plume sparso --- SRC109, V1, Q1/2: l'agente finisce a 2532\,m (partenza 2004\,m).
           Il plume diffuso annulla il gradiente locale: l'agente degenera in random walk
           e nei video appare quasi fermo per tutta la durata dell'episodio.}
\end{figure}

\vspace{1.2cm}

\begin{figure}[H]
  \centering
  \includegraphics[width=0.62\linewidth]{../../evaluations/evaluations_FCM/fcm_pure_50m/SRC127/V2/ep01_chunk2_trajectory.png}
  \caption{Massimo locale --- SRC127, V2, Q3/4: l'agente percorre 886\,m verso la sorgente
           (da 1008\,m a 122\,m) ma si blocca su un picco locale di concentrazione
           e non riesce a uscirne fino al timeout.}
\end{figure}

% ══════════════════════════════════════════════════════════════════════════════
\section{Confronto FCM vs.\ RL}

\begin{table}[H]
\centering
\renewcommand{\arraystretch}{1.3}
\small
\begin{tabular}{L{4cm} C{3.5cm} C{3.5cm}}
\toprule
\rowcolor{headerbg}
\textcolor{white}{\textbf{Metrica}} &
\textcolor{white}{\textbf{FCM ($r_s{=}50\,\text{m}$)}} &
\textcolor{white}{\textbf{RL (MaskablePPO)}} \\
\midrule
\rowcolor{rowalt}
SR globale          & 32{,}9\%  & \textbf{96{,}9\%} {\scriptsize\textcolor{green!50!black}{(+64{,}0\%)}} \\
Step medi (succ.)   & 182{,}7 ($\sim$30{,}5\,min)  & \textbf{129{,}2 ($\sim$21{,}5\,min)} \\
\rowcolor{rowalt}
V0                  & 39{,}7\%  & \textbf{100{,}0\%} {\scriptsize\textcolor{green!50!black}{(+60{,}3\%)}} \\
V1                  & 16{,}0\%  & \textbf{99{,}0\%}  {\scriptsize\textcolor{green!50!black}{(+83{,}0\%)}} \\
\rowcolor{rowalt}
V2                  & 23{,}7\%  & \textbf{88{,}7\%}  {\scriptsize\textcolor{green!50!black}{(+65{,}0\%)}} \\
V3                  & 51{,}9\%  & \textbf{99{,}7\%}  {\scriptsize\textcolor{green!50!black}{(+47{,}8\%)}} \\
\rowcolor{rowalt}
Q1/4                & 50{,}5\%  & \textbf{100{,}0\%} {\scriptsize\textcolor{green!50!black}{(+49{,}5\%)}} \\
Q1/2                & 11{,}5\%  & \textbf{93{,}5\%}  {\scriptsize\textcolor{green!50!black}{(+82{,}0\%)}} \\
\rowcolor{rowalt}
Q3/4                & 36{,}5\%  & \textbf{97{,}1\%}  {\scriptsize\textcolor{green!50!black}{(+60{,}6\%)}} \\
\bottomrule
\end{tabular}
\caption{Confronto diretto FCM ottimale vs.\ modello RL (\texttt{ppo\_20260516\_143937}, 1560 episodi).
         Stesse sorgenti di test, stesso ambiente, stessa soglia di successo.}
\end{table}

% ══════════════════════════════════════════════════════════════════════════════
\section{Discussione e Conclusione}

\subsection{Limiti strutturali e casi critici}

Il Field Climbing Method \`e un metodo \textbf{puramente locale e reattivo}:
non mantiene memoria dell'episodio, non sfrutta informazioni di vento e corrente,
e non dispone di alcun meccanismo di esplorazione globale.
Questo comporta due categorie di fallimento sistematico:

\begin{enumerate}
    \item \textbf{Campo localmente piatto.}
    In zone lontane dalla sorgente, o in scenari di vento debole (V1), il plume
    \`e cos\`i diffuso che tutti i sensori restituiscono concentrazioni quasi identiche.
    Il gradiente stimato \`e nullo o dominato dal rumore: l'agente degenera in un
    random walk con passi da 10\,m e accumula spostamenti confinati a poche
    decine di metri su scala di centinaia di metri di distanza dalla sorgente.
    V1 \`e il caso pi\`u critico: success rate del 16{,}0\% contro il 99{,}0\% del modello RL
    (+83pp), confermando che il plume diffuso non \`e un problema di parametrizzazione
    del sensor range ma una limitazione intrinseca dell'approccio gradient-based.

    \item \textbf{Gradiente locale ingannevole.}
    In un ambiente costiero con correnti complesse, la struttura del plume pu\`o
    presentare massimi locali di concentrazione distanti dalla sorgente reale.
    FCM segue il gradiente locale senza meccanismi di esplorazione forzata:
    viene attratto da questi massimi locali e si blocca.
    Il chunk Q1/2 registra le performance peggiori per tutti i sensor range
    (11{,}5\% nel caso migliore), suggerendo che il campo a met\`a simulazione
    abbia una struttura di plume pi\`u complessa o meno direzionata.
\end{enumerate}

\subsection{Conclusione}

Il gap di 64 punti percentuali tra FCM (32{,}9\%) e RL (96{,}9\%) quantifica
il contributo dell'apprendimento per rinforzo in questo problema.
Il modello RL ha appreso strategie non riducibili al gradiente locale: sfruttamento
delle informazioni di vento e corrente, navigazione in zone a bassa concentrazione,
recupero da situazioni di stallo.
Queste strategie emergono dall'ottimizzazione su milioni di episodi con scenari
diversificati e non sono accessibili a un metodo che decide step-by-step
senza contesto storico.

\end{document}
